\label{sec:1}

\section{Motivation}
\label{sec:1.1}
Pursuit-evasion problems model scenarios where agents compete to intercept or evade one another. They are a well-studied problem in aerospace systems research, with practical applications in autonomous air combat, unmanned aerial vehicle (UAV) swarm coordination, and airspace collision avoidance.

Approaching real-time solutions to these problems is challenging: the state space of even a simple scenario is continuous and high-dimensional. Exact solutions are computationally intractable\cite{mcgrew2010}, but Markov Decision Processes provide a performant framework for efficient approximation\cite{bertram2d}, and have been shown to scale to multi-agent teams in continuous 3D state spaces\cite{bertram3d}.

This project implements and extends the FastMDP algorithm of Bertram et al.\cite{bertram2d} in a 3D kinematic simulation of fixed-wing aircraft. The primary focus is pursuit-evasion behaviour, investigated through novel reward functions and performance analysis. A key contribution is demonstrating that the same simulation and decision-making framework generalises to other problem spaces entirely, including collaborative flight and search-and-rescue, with changes only to the reward structure.

The project is implemented in Python 3.13.2, using \texttt{numpy} and \texttt{numba} for numerical computation.

\section{Project Aims}
\label{sec:1.2}
The core aims of this project are to:
\begin{enumerate}[noitemsep]
    \item Implement a kinematically accurate 3D simulation of fixed-wing aircraft using a pseudo-6DoF model.
    \item Implement the FastMDP algorithm to drive behaviours between competing agents.
    \item Design reward functions that produce valid pursuit-evasion behaviour.
    \item Analyse the engineering performance of the simulation and decision-making framework.
\end{enumerate}

After completing the core project, the following extensions were investigated:

\begin{enumerate}[resume,noitemsep]
    \item Scale the simulation from 1v1 encounters to multi-agent teams.
    \item Demonstrate that the decision-making framework is applicable to alternative problem spaces, requiring changes only to the reward structure.
\end{enumerate}

\section{Success Criteria}
\label{sec:1.3}
The following criteria were used to evaluate the success of the project:
\begin{enumerate}
    \item \textbf{Behavioural validity}. The implemented FastMDP algorithm produces valid pursuit-evasion behaviour. This is the primary functional goal of the project and directly validates the correctness of the decision-making and reward framework.
    \item \textbf{Kinematic validity}. Agent kinematics remain within realistic bounds across all simulation runs, including velocity, turn rate and altitude constraints. This validates the correctness of the kinematic model independently of the decision-making framework.
    \item \textbf{Computational performance}. Agent actions are computed within a timeframe suitable for real-time simulation with appropriate optimisations investigated and evaluated.
    \item \textbf{Extension validity}. Multi-agent teams and alternative problem spaces produce valid and consistent behaviour under the same framework.
\end{enumerate}

The remainder of the dissertation covers the preparation, implementation and evaluation carried out to meet these criteria.